% TOP_PROBLEM: {"id": 2304022, "title": "Research Problems in Function Theory — Problem 4.22", "category_id": 8, "category": {"id": 8, "name": "analysis", "display_name": "Analysis", "description": "Limits, continuity, calculus, and function theory.", "slug": "analysis", "order_index": 8, "created_at": "2026-07-31T15:26:25.671Z"}, "status": "open"}
% FIRST_POSTED: null
% ENTRY_KIND: statement_only
% Imported from: batch13.tex; UnsolvedMath ID 2304022
\documentclass[11pt]{book}
\usepackage{fontspec}
\setmainfont{Latin Modern Roman}
\setsansfont{Latin Modern Sans}
\usepackage{amsmath,amssymb,mathtools}
\usepackage{unicode-math}
\setmathfont{Latin Modern Math}
\usepackage{xurl}
\usepackage[hidelinks]{hyperref}
\newcommand{\sref}[2]{\hyperlink{source:#1:#2}{[S#2]}}
\newcommand{\cataloguescope}[1]{\paragraph{Mathematical question.}#1}
\begin{document}
% UnsolvedMath ID 2304022; source catalogue code AMR-022-4022
\setcounter{section}{2304021}
\section{Research Problems in Function Theory — Problem 4.22}\label{2304022}
{\small\sffamily UnsolvedMath ID 2304022 · Analysis}
\subsection{Definitions and mathematical statement}

\paragraph{Shared notation.}$\mathbb N=\{1,2,\ldots\}$, and $\mathbb Z,\mathbb Q,\mathbb R,\mathbb C$ denote the integers, rationals, reals and complex numbers. Any other conventions are specified locally.


Let $n\ge1$ and let $f(z)=\prod_{j=1}^n(z-z_j)$ be a monic polynomial of degree $n$, with zeros repeated according to multiplicity. Its filled unit lemniscate is the compact set
\[E_f=\{z\in\mathbb C:|f(z)|\le1\}.
\]
The monic normalization is the one used in the source's Problem 4.7 and its following lemniscate questions.Assume $z_1,\ldots,z_n\in\overline{\mathbb D}$, so $0\in E_f$. Let
\[L_f=\inf\{\ell(\gamma):\gamma:[0,1]\to E_f\text{ is a rectifiable continuous path},\ \gamma(0)=0,\ |\gamma(1)|=1\},\]
with infimum $+\infty$ if no rectifiable path exists. Define the worst-case quantity
\[L(n)=\sup_{z_1,\ldots,z_n\in\overline{\mathbb D}}L_f.
\]
The source records existence of a path from $0$ to $\mathbb T$ in $E_f$.
\paragraph{Statement recorded in the supplied source.}
Estimate the shortest-path length $L_f$ and determine optimal degree-dependent upper bounds, equivalently the size and growth of $L(n)$. In particular, is the worst-case length unbounded as $n\to\infty$, and how slowly can the optimal upper bound grow? \sref{2304022}{2}
\subsection{Short English statement}
How long may the shortest path from the origin to the unit circle inside a polynomial’s lemniscate have to be, when all zeros lie in the unit disc?
\subsection{Sources}
\begin{description}
\item[\hypertarget{source:2304022:1}{S1}] ULAM.AI and the UnsolvedMath Contributors, \emph{UnsolvedMath: A Curated Collection of Open Mathematics Problems}, record 2304022 (AMR-022-4022). CC BY 4.0. \url{https://huggingface.co/datasets/ulamai/UnsolvedMath}.
\item[\hypertarget{source:2304022:2}{S2}] Walter K. Hayman and Edward F. Lingham, \emph{Research Problems in Function Theory}, arXiv:1809.07200, 2018. Problem 4.22, its complete Update 4.22, and the relevant chapter notation. \url{https://arxiv.org/pdf/1809.07200}.
\end{description}
\label{end:2304022}
\end{document}
